% LaTeX Template For MATH 490 @ VCU
\documentclass{article}

\usepackage{hyperref}
\usepackage{amsmath}
\usepackage{amsthm}
\usepackage{amssymb}
\usepackage{geometry}
\geometry{a4paper, portrait, margin=2cm}


\begin{document}

\title{Project Proposal}

\author{Tom Barnes}

\maketitle              % typeset the title of the contribution

\subsection{Introduction}
For this project I am going to focus on a result from elementary number theory
as I am accustomed to working with natural numbers and inductive proofs in lean. 
I will aim to formalise the existence part of the fundamental theorem of arithmetic, 
i.e. for any natural number $n > 1$, there exists prime numbers $p_1, p_2, \dots ,p_k$ such
that $n=p_1p_2\cdots p_k$. If I have time, or if I complete all the lemmas and the main result
in under 200 lines, I might try to also prove that the primes are unique up to reordering.

\subsection{Proof Outline}
The basic proof I will formalise will be a proof by induction as follows:
\begin{itemize}
  \item If $n>1$ is prime, then its factorisation is itself ($n=2$ is the base case)
  \item If $n>2$ is composite then it has a factorisation $n=ab$ with $1<a,b<n$ 
  \item Apply the induction hypothesis to get $a=p_1 \cdots p_k$ and $b = q_1 \cdots q_l$
  \item Combine the factorisations of $a$ and $b$ to get a factorisation for $n$
\end{itemize}

\subsection{Definitions}
Here are the basic definitions I will use:
\begin{itemize}
  \item I will use Mathlib's definition for primality
  \item I will use my own definition of a prime factorisation of $n$ as simply a list of 
  natural numbers that are primes and also whose product is equal to $n$
\end{itemize}

\subsection{Results}
Here are the lemmas I intend to prove: (more will probably be required)
\begin{enumerate}
  \item Every prime number has a prime factorisation (a singleton list containing itself)
  \item An empty list represents a prime factorisation of one
  \item Appending two prime factorisations gives another prime factorisation 
  \item Every natural number greater than one is either prime or has a factorisation into 
  two natural numbers larger than one and smaller than itself
  \item Strong induction for the main result
\end{enumerate} 

\subsection{Mathlib Usage}
As mentioned earlier, I will use the primality predicate from Mathlib as a definition. I will 
also use Mathlib's definition of natural numbers and lean's builtin linked list data structure.
I will also use results from Mathlib relating to basic arithmetic, divisibility and operations
on lists. I will not be using any functions from Mathlib relating to factorisation or prime 
numbers, they will be formalised explicitly when required. I will likely use the standard induction
tactic seen in the Natural Number game.

\end{document}
